% !TeX program = lualatex
% !BIB program = biber
\documentclass[11pt,a4paper]{article}

\usepackage[a4paper,margin=2.5cm]{geometry}
\usepackage{fontspec}
\IfFontExistsTF{Times New Roman}{
  \setmainfont{Times New Roman}
}{
  \setmainfont{TeX Gyre Termes}
}
\setmonofont{TeX Gyre Cursor}

\usepackage[ngerman]{babel}
\usepackage{csquotes}
\usepackage{setspace}
\onehalfspacing
\usepackage{microtype}
\usepackage{graphicx}
\graphicspath{{abbildungen/}}
\usepackage{tikz}
\usetikzlibrary{arrows.meta,positioning}
\usepackage{booktabs}
\usepackage{tabularx}
\usepackage{enumitem}
\usepackage{titlesec}
\usepackage{xcolor}
\usepackage{listings}
\usepackage[font=small,labelfont=bf]{caption}
\usepackage{fancyhdr}
\usepackage{hyperref}
\usepackage[nameinlink,noabbrev]{cleveref}
\crefname{figure}{Abbildung}{Abbildungen}
\Crefname{figure}{Abbildung}{Abbildungen}
\crefname{table}{Tabelle}{Tabellen}
\Crefname{table}{Tabelle}{Tabellen}
\crefname{lstlisting}{Listing}{Listings}
\Crefname{lstlisting}{Listing}{Listings}
\crefname{section}{Abschnitt}{Abschnitte}
\Crefname{section}{Abschnitt}{Abschnitte}
\crefname{subsection}{Abschnitt}{Abschnitte}
\Crefname{subsection}{Abschnitt}{Abschnitte}
\usepackage[
  backend=biber,
  style=authoryear,
  sorting=nyt,
  maxcitenames=2,
  maxbibnames=99,
  doi=true,
  isbn=true,
  url=false,
  eprint=false
]{biblatex}
\addbibresource{literatur.bib}

% Zentrale Projektdaten: nur hier ausfuellen/aendern.
\newcommand{\Autor}{Frithjof Beims}
\newcommand{\Matrikelnummer}{2765521}
\newcommand{\Dozent}{Prof. Dr. Thomas Clemen}
\newcommand{\Praktikumsgruppe}{StudyMummy}
\newcommand{\Projektname}{StudyMummy}
\newcommand{\Arbeitstitel}{Kognitive Architektur eines agentischen Systems}
\newcommand{\RepositoryURL}{https://github.com/Fridi286/StudyMummy}
\newcommand{\RepositoryAnzeige}{github.com/Fridi286/StudyMummy}
\newcommand{\Abgabedatum}{16. August 2026}
\newcommand{\Vertiefungskapitel}{Kap. 02 -- Kognitive Architektur}


% Fussnoten werden in 10 pt gesetzt.
\renewcommand{\footnotesize}{\fontsize{10pt}{12pt}\selectfont}
\setlength{\parindent}{0pt}
\setlength{\parskip}{0pt}
\setlength{\emergencystretch}{1.5em}
\setlist{itemsep=0.1em,topsep=0.2em,parsep=0pt,partopsep=0pt}
\titlespacing*{\section}{0pt}{2.0ex plus 0.5ex minus 0.2ex}{0.6ex}
\titlespacing*{\subsection}{0pt}{1.5ex plus 0.3ex minus 0.2ex}{0.4ex}
\titlespacing*{\subsubsection}{0pt}{1.2ex plus 0.3ex minus 0.2ex}{0.3ex}
\setcounter{secnumdepth}{3}
\setcounter{tocdepth}{2}

\definecolor{hawblue}{HTML}{005AA9}
\hypersetup{
  colorlinks=true,
  linkcolor=black,
  citecolor=hawblue,
  urlcolor=hawblue,
  pdftitle={\Arbeitstitel},
  pdfauthor={\Autor}
}

\pagestyle{fancy}
\fancyhf{}
\fancyhead[L]{WP Agentic AI}
\fancyhead[R]{\Autor}
\fancyfoot[C]{\thepage}
\renewcommand{\headrulewidth}{0.4pt}
\setlength{\headheight}{14pt}

\lstset{
  basicstyle=\ttfamily\small,
  numbers=left,
  numberstyle=\scriptsize\color{gray},
  frame=single,
  breaklines=true,
  columns=fullflexible,
  captionpos=b,
  showstringspaces=false
}

\begin{document}

\hypersetup{pageanchor=false}
\begin{titlepage}
  \thispagestyle{empty}
  \begin{center}
    {\large Hochschule für Angewandte Wissenschaften Hamburg\par}
    \vspace{0.4cm}
    {\large Bachelor Informatik -- Sommersemester 2026\par}
    \vfill
    {\Large Hausarbeit im Modul\par}
    \vspace{0.25cm}
    {\LARGE\bfseries WP Agentic AI\par}
    \vspace{0.8cm}
    {\huge\bfseries \Arbeitstitel\par}
    \vspace{0.5cm}
    {\Large Analyse des Praktikumsprojekts \Projektname\par}
    \vfill
    \begin{tabular}{@{}ll@{}}
      Autor: & \Autor \\
      Matrikelnummer: & \Matrikelnummer \\
      Dozent: & \Dozent \\
      Praktikumsgruppe: & \Praktikumsgruppe \\
      Projekt: & \Projektname \\
      GitHub-Repository: & \href{\RepositoryURL}{\RepositoryAnzeige} \\
      Vertiefungskapitel: & \Vertiefungskapitel \\
      Abgabedatum: & \Abgabedatum
    \end{tabular}
    \vfill
  \end{center}
\end{titlepage}

% Arabische Nummerierung beginnt auf der ersten Seite nach dem Deckblatt.
\clearpage
\hypersetup{pageanchor=true}
\pagenumbering{arabic}
\setcounter{page}{1}
\begingroup
\singlespacing
\small
\tableofcontents
\endgroup

\setcounter{section}{-1}
\section{Systemüberblick}
\label{sec:systemueberblick}

StudyMummy ist eine webbasierte Lernplattform, die hochgeladene Lernunterlagen
in Aufgaben, Quizfragen und Merkblätter überführt und Lernende in einem
sokratischen Chat adaptiv unterstützt. Der agentische Kern verarbeitet pro Chat-Turn eine
authentifizierte Nachricht zusammen mit Aufgabe, Hilfestufe, Dialoghistorie und
abgerufenem Dokumentwissen. Planner, Tutor und Reviewer besitzen getrennte
Ziele, Befugnisse und laufbezogene Zustände. Über typisierte Nachrichten kann
der Reviewer einen Entwurf freigeben, eine Tutorrevision verlangen oder eine
begrenzte Neuplanung anstoßen. Der Orchestrator verteilt diese Nachrichten und
begrenzt den Ablauf, trifft aber keine fachliche Entscheidung.

Technisch besteht das System aus einer Angular-22-Single-Page-Application,
einer asynchronen FastAPI-Anwendung mit Pydantic und SQLAlchemy sowie
PostgreSQL~15 mit \texttt{pgvector}. Ein OpenAI-kompatibler Adapter stellt
strukturierte Aufrufe großer Sprachmodelle (Large Language Models, LLMs) und
Function Calling bereit. Docker Compose trennt Entwicklungs- und
Produktionsprofile. Im Produktivprofil liefert Nginx das
Frontend aus und leitet API- und WebSocket-Verkehr an das Backend weiter.
Persistiert werden unter anderem Nutzer, Dokumente, Vektorchunks, Aufgaben,
Sitzungen, Chatlogs, Lernprofile und Gamification-Zustände. Die zentrale Frage
lautet: \emph{\Forschungsfrage}

\section{Agent-Analyse (Kap. 01)}
\label{sec:agent-analyse}

\subsection{PEAS-Analyse}

Ein Agent ist durch seine Wahrnehmung und Wirkung in einer Umgebung bestimmt,
PEAS zwingt dabei zur expliziten Operationalisierung des Erfolgs
\parencite[Kap.~2]{russellNorvig2021}. Für den Tutor-Workflow ergibt sich
\cref{tab:peas}.

\begin{table}[htbp]
  \centering
  \small
  \begin{tabularx}{\linewidth}{@{}lX@{}}
    \toprule
    PEAS-Komponente & Projektspezifische Ausprägung und Codebeleg \\
    \midrule
    Performance & Fachlich plausible, nicht leere und zur Hilfestufe passende
    Antwort, Erfüllung von \texttt{success\_criteria}, korrekte Werkzeugwahl
    und Reviewer-Freigabe. Eine einzige numerische Nutzenfunktion existiert
    nicht. Die separaten Evaluationsmetriken prüfen sokratische Führung,
    Vermeidung direkter Lösungen, Trace-ID und Tool-Abdeckung. \\
    Environment & Authentifizierte Lernende, Angular-Oberfläche, aktuelle
    Aufgaben und Dokumente, PostgreSQL-Zustand, LLM-Anbieter sowie Resultate
    interner Werkzeuge. \\
    Actuators & Tutor-Text als primäre Wirkung, zusätzlich die im Chat
    tatsächlich freigebbaren Werkzeuge \texttt{evaluate\_answer},
    \texttt{award\_coins\_and\_exp} und \texttt{add\_calendar\_note}. Die
    letzten beiden verändern persistenten Zustand. \\
    Sensors & Nachricht, Zeit, Dialoghistorie, aktive Aufgabe, Hilfestufe,
    Frontend-Kontext, bis zu drei semantisch ähnliche Dokumentchunks und
    Werkzeugergebnisse. Die Bildung des \texttt{AgentContext} erfolgt im
    Endpunkt \path{backend/app/api/v1/endpoints/agent.py}. \\
    \bottomrule
  \end{tabularx}
  \caption{PEAS-Analyse des agentischen Tutor-Workflows.}
  \label{tab:peas}
\end{table}

\subsection{Taxonomie und Umgebungsdimensionen}

StudyMummy ist am treffendsten ein \emph{begrenzt modellbasierter,
zielorientierter und deliberativer} Softwareagent. Modellbasiert ist er, weil
Sitzung, Aufgabe, Hilfestufe und Historie einen internen Zustand bilden.
Zielorientiert ist er durch das je Turn erzeugte \texttt{objective} und ein bis
drei Erfolgskriterien. Der JSON-Plan und die nachfolgende Ausführung gehen über
einen Reflexagenten hinaus. Eine Einordnung als Utility-Agent wäre dagegen zu
stark: Weder werden konkurrierende Pläne bewertet noch wird erwarteter Nutzen
maximiert. Auch ein Lernagent im klassischen Sinn ist nicht vollständig
implementiert. Zwar existieren Lernprofil und Update-Werkzeug, doch die
Capability-Policy gibt dieses Werkzeug im Chat nie frei und das Profil fließt
nicht in den \texttt{AgentContext} ein.

Die Umgebung ist nur \emph{partiell beobachtbar}, da Wissen, Motivation und
tatsächliches Verständnis der lernenden Person indirekt aus Text erschlossen
werden. Sie ist \emph{sequenziell} und \emph{dynamisch}: Frühere Antworten,
Aufgabenstatus und Dokumente beeinflussen spätere Entscheidungen und können
sich zwischen Turns ändern. Datenbankoperationen sind überwiegend
deterministisch, LLM-Ausgaben und Nutzerverhalten jedoch stochastisch. Aktionen
und Zustände sind in der Implementierung diskret typisiert, der relevante
Zustandsraum bleibt wegen freier Sprache sehr groß. Die eigenen
Werkzeugtransitionen sind bekannt, die Reaktion von Mensch und Modell nur
teilweise. Innerhalb des Systems interagieren mehrere logische Agentenrollen,
gegenüber der Lernumgebung ist der Workflow zugleich ein gemeinsamer
Tutorakteur.

\section{Kognitive Architektur (Kap. 02) -- Vertiefungskapitel}
\label{sec:kognitive-architektur}

\subsection{Ist-Architektur: Perception, Memory, Planning und Action}

Als Beobachtungsrahmen wird zwischen Perception, Memory, Planning und Action
unterschieden. Diese Zuordnung ist eine konzeptionelle Analyse des Codes, keine
Behauptung, dass vier eigenständige Laufzeitmodule existieren. CoALA präzisiert
für Sprachagenten insbesondere explizites Arbeits- und Langzeitgedächtnis,
interne und externe Aktionen sowie einen wiederholten Entscheidungszyklus
\parencite{sumersEtAl2024}. StudyMummy realisiert Teile davon als
FastAPI-Endpunkt, Datenbankservices und drei LLM-Rollen.

Der konkrete Kontrollfluss beginnt in \texttt{chat()}: Eingabe filtern,
Eigentum an Aufgabe und Dokument prüfen, Sitzungszustand laden und Hilfestufe
aktualisieren, Nutzerturn speichern, RAG-Kontext abrufen und den typisierten
\texttt{AgentContext} bilden. Der Orchestrator ruft danach Planner, Tutor und
optional Reviewer auf. Erst außerhalb des Orchestrators speichert der Endpunkt
die finale Antwort und ergänzt den Trace-Schritt \enquote{remember}. Dieser
letzte Punkt ist relevant: Persistenz ist vorhanden, aber nicht Teil einer
lernfähigen Rückkopplung des Orchestrators.

\subsubsection{Perception}

Die Rohwahrnehmung besteht aus freier Nutzereingabe und Datenbankzustand.
\texttt{filter\_user\_input()} normalisiert Unicode, entfernt unsichtbare
Steuerzeichen, reduziert Leerraum, begrenzt auf 4.000 Zeichen und blockiert
einige reguläre Ausdrücke für direkte Prompt-Injection. Anschließend verknüpft
\texttt{\_resolve\_study\_context()} die Nachricht nur mit Aufgaben und
Dokumenten des authentifizierten Nutzers. Die Hilfestufe wird anhand weniger
Schlüsselwörter deterministisch auf 1 bis 4 angehoben. Für RAG wird die Frage
mit dem vertrauenswürdigen Aufgabenkontext kombiniert, der Abruf filtert nach
\texttt{user\_id} und optional \texttt{document\_id}.

Die resultierende Repräsentation ist zweckmäßig, aber selektiv. Der Planner
erhält weder Dokumentinhalt noch Dialoghistorie, sondern nur Nachricht,
Hilfestufe, Aufgabe und ein boolesches Kontextmerkmal. Das reduziert
Injection-Risiken, verhindert aber dokumentabhängige Planung. Umgekehrt kann
\texttt{extra\_context} bis zu 12.000 Zeichen aus dem Frontend enthalten und
wird nicht durch den Eingabefilter geprüft. RAG-Text wird im Tutor-Prompt zwar
als nicht vertrauenswürdig markiert, eine strukturelle Trennung von Daten und
Anweisungen findet jedoch nicht statt. Fehlgeschlagener Retrieval-Zugriff führt
laut \texttt{rag\_service.py} still zu leerem Kontext und ist für Planner und
Nutzer nicht unterscheidbar von einem legitimen Treffer ohne Ergebnis.

\subsubsection{Memory}

Das Arbeitsgedächtnis ist explizit in \texttt{Session} und
\texttt{WorkingMemory} modelliert: aktive Aufgabe, Hilfestufe und kompletter
Dialog werden vor jedem Turn geladen. In CoALA ist ein solches
Arbeitsgedächtnis eine über LLM-Aufrufe persistente Datenstruktur und nicht nur
das Kontextfenster \parencite{sumersEtAl2024}. StudyMummy erfüllt dies teilweise,
kopiert jedoch die gesamte Historie ohne Tokenbudget, Zusammenfassung oder
Vergessen in jeden Tutoraufruf. Lange Sitzungen wachsen daher unbeschränkt bis
zum Providerlimit.

Episodisch persistieren \texttt{ChatLog}-Datensätze Rolle, Inhalt, Zeitpunkt
und die finale Aktion. Pläne, Reviewer-Feedback, Werkzeugargumente,
Werkzeugresultate und Erfolg oder Fehlschlag werden nicht gespeichert und
können später nicht gezielt abgerufen werden. Semantisches Gedächtnis besteht
aus hochgeladenen Dokumenten: Der Analyzer zerlegt Volltext mit
\texttt{textwrap.wrap} in ungefähr 1.000 Zeichen lange Chunks, erzeugt
Embeddings und speichert 1.536-dimensionale Vektoren in PostgreSQL. Online
werden die drei Chunks mit kleinster Kosinusdistanz geladen. Eine
Ähnlichkeitsschwelle, ein Reranker, überlappende oder strukturerhaltende Chunks
und Quellenzitate fehlen. Lernprofile und Confidence Scores bilden zwar einen
weiteren Langzeitspeicher, bleiben im analysierten Chatpfad aber ungenutzt.

\subsubsection{Planning}

Der Planning Agent erzeugt mit Temperatur 0 ein \texttt{AgentPlan}-Objekt aus
Intent, genau einer Aktion, Ziel, kurzer Entscheidungsbasis, Werkzeugnamen und
Erfolgskriterien. Pydantic validiert das JSON. Bei Provider- oder
Strukturfehlern greift ein deterministischer, schlüsselwortbasierter Fallback.
Danach schneidet \texttt{SAFE\_TOOL\_POLICY} die vorgeschlagenen Rechte auf ein
aktionsabhängiges Set zurück. Diese Kombination ist nachvollziehbarer als ein
freier Prompt und implementiert eine begrenzte deliberative Auswahl.

Sie ist jedoch weder hierarchische Zielzerlegung noch Closed-Loop-Planung. Der
Planner kennt Werkzeugbeobachtungen nicht, bewertet keine Alternativen und wird
nach Toolfehler oder Reviewer-Ablehnung nicht erneut aufgerufen. Der Tutor darf
innerhalb einer Aktion bis zu fünf Function-Calling-Runden ausführen, diese
lokale Schleife ähnelt ReAct, während der übergeordnete Plan unverändert bleibt.
ReAct begründet gerade den Nutzen, Reasoning und externe Beobachtung
abwechselnd zur Planaktualisierung zu verwenden \parencite{yaoEtAl2023react}.
StudyMummy nutzt die Beobachtung nur im Tutor-Dialog, nicht im typisierten
Planungszustand.

\subsubsection{Action}

Action umfasst Textausgabe und echte Function Calls. Der LLM-Adapter bietet nur
die freigegebenen JSON-Schemas an, blockiert unbekannte Namen erneut, parst
Argumente und gibt Werkzeugresultate als \texttt{tool}-Nachricht an den Tutor
zurück. JSON-Fehler und Python-Ausnahmen werden in eine Fehlerbeobachtung
umgewandelt. Praktisch erreichbar sind Antwortbewertung, Vergabe virtueller
Coins/Erfahrung und ein Kalendereintrag. Die Belohnung bindet die Identität über
einen Request-Kontext, begrenzt den Betrag, prüft Aufgabeneigentum unter
Datenbanksperre und verhindert Doppelbelohnung.

Der Reviewer sieht anschließend nur Plan und fertigen Entwurf. Er kann Text
ersetzen, aber einen bereits angelegten Termin oder eine Belohnung nicht
zurücknehmen. Bei ungültigem Reviewer-JSON wird der Entwurf sogar freigegeben.
Außerdem verlangt der Tutor-Prompt ein Lernprofil-Update, obwohl dieses Tool
laut Policy nie im Scope liegt. Diese Diskrepanz zwischen prozeduralem Wissen
(Prompt) und ausführbarem Aktionsraum ist ein konkreter Architekturfehler.

\subsection{Kritische Analyse und Grenzen}

Die Schichten sind grundsätzlich erkennbar, aber ihre Rückkopplung ist schwach.
Erstens ist die persistierte Historie eher ein Gesprächsarchiv als episodisches
Gedächtnis: Es gibt weder selektiven Abruf noch Konsolidierung und Reflexion.
Generative Agents zeigen ein Gegenmodell, in dem Erfahrungen gespeichert,
dynamisch abgerufen und zu höherwertigen Reflexionen verdichtet werden, die
Ablation der Autoren weist Wahrnehmung, Planung und Reflexion jeweils als
relevante Beiträge aus \parencite{parkEtAl2023}. Eine Übertragung muss dennoch
domänenspezifisch evaluiert werden, weil glaubwürdiges Simulationsverhalten
nicht mit Lernerfolg gleichzusetzen ist.

Zweitens ist Kontrolle überwiegend vor der Aktion (Tool-Whitelist) oder nach
der Antwort (Reviewer) platziert, nicht zwischen Beobachtung und irreversibler
Wirkung. Das System speichert auch kein überprüftes Outcome, aus dem spätere
Entscheidungen lernen könnten. Reflexion nutzt verbales Feedback in einem
episodischen Puffer, um Folgeversuche zu verbessern
\parencite{shinnEtAl2023reflexion}, zugleich hängt dieses Verfahren von der
Qualität der Selbstbewertung ab und garantiert keinen Erfolg. Für StudyMummy
folgt daraus nicht, ungeprüfte Selbstkritik zu akkumulieren, sondern nur
validiertes Reviewer- und Werkzeugfeedback strukturiert zurückzuführen.

\subsection{Verbesserungsvorschlag 1: Begrenztes, konsolidiertes Episodengedächtnis}
\label{sec:verbesserung-eins}

Technisch sollte eine Tabelle \texttt{agent\_episodes} ergänzt werden, deren
Datensatz \texttt{user\_id}, \texttt{session\_id}, \texttt{task\_id}, den
validierten Plan, ausgeführte Aktion, Werkzeugstatus, Reviewer-Urteil, eine
knappe Ergebniszusammenfassung, Zeit und optional ein Embedding enthält. Ein
deterministischer \texttt{MemoryWriter} schreibt erst nach Review und Commit,
interne Gedankengänge werden nicht gespeichert. Ein
\texttt{Memory\allowbreak{}Retriever}
lädt vor der Planung höchstens drei nutzer- und aufgabengefilterte Episoden nach
Semantik, Aktualität und Ergebnisstatus. Parallel fasst ein
\texttt{session\_summary} ältere Chatturns ab einem festen Tokenbudget zusammen,
die letzten Turns bleiben wörtlich. Lernprofil und relevante Confidence Scores
werden als eigene, vertrauenswürdige Felder in \texttt{AgentContext} gebunden.

Dies überträgt CoALAs Trennung von Arbeits-, episodischem und semantischem
Gedächtnis sowie den Abruf früherer Episoden in die vorhandene SQLAlchemy- und
pgvector-Architektur \parencite{sumersEtAl2024}. Erwartet werden stabilere lange
Sitzungen, gezieltere Wiederaufnahme früherer Fehler und begrenzte Promptgröße.
Voraussetzungen sind ein Lösch- und Aufbewahrungskonzept, Migrationen und ein
Evaluationssatz mit Mehrsitzungsverläufen. Zusätzliche Embeddings und Abrufe
erhöhen Latenz und Kosten, Schema, Zusammenfasser und Retrieval-Scoring erhöhen
Komplexität und Wartungsaufwand. Die Sicherheitsoberfläche wächst durch
\emph{Memory Poisoning}: manipulierte Inhalte könnten dauerhaft wiederkehren.
Daher sind Provenienz, Nutzerisolation, TTL, Löschbarkeit und das Speichern nur
validierter Outcomes zwingend. Auch Reflexionen dürfen nicht ungeprüft als
Fakten behandelt werden.

\subsection{Verbesserungsvorschlag 2: Typisierte Plan--Act--Observe--Replan-Schleife}
\label{sec:verbesserung-zwei}

Der Orchestrator sollte ein explizites Objekt \texttt{AgentState} führen. Es
enthält \texttt{plan}, \texttt{observations}, \texttt{attempt},
\texttt{pending\_\allowbreak side\_\allowbreak effect}, \texttt{review} und
\texttt{status}. Der Planner ergänzt pro Aktion eine erwartete Beobachtung und
eine Fehlerstrategie. Zunächst werden nur lesende Werkzeuge ausgeführt, deren
Resultate werden gegen Pydantic-Schemas validiert. Bei Toolfehler oder
Reviewer-Ablehnung erhält der Planner die kompakte Beobachtung und darf höchstens
einmal neu planen. Schreibende Tools werden als \texttt{pending} materialisiert
und erst nach Policy-Prüfung sowie -- beim Kalender -- Nutzerbestätigung
committet. Ein Datenbank-Checkpoint vor der Wirkung ermöglicht sauberen Abbruch,
harte Grenzen für Schritte, Zeit und Toolbudget verhindern Schleifen.

Die Ausgestaltung setzt ReActs alternierende Kopplung von Planen und
Umweltbeobachtung um, ohne Chain-of-Thought offenzulegen
\parencite{yaoEtAl2023react}. Validiertes Reviewer-Feedback kann zusätzlich als
sprachliches Fehlersignal für den nächsten Versuch dienen, entsprechend der
Grundidee von Reflexion \parencite{shinnEtAl2023reflexion}. Der erwartete Nutzen
liegt in kontrollierter Fehlerbehandlung, überprüfbaren Abbruchbedingungen und
dem Schließen der heutigen Lücke, in der Review erst nach Seiteneffekten
erfolgt. Dem stehen mindestens ein weiterer LLM-Aufruf im Fehlerfall, höhere
Latenz und Kosten, komplexere Transaktionen und mehr Zustände im Test gegenüber.
Jeder zusätzliche Rückkopplungskanal kann zudem Injection weitertragen, nur
strukturierte Beobachtungen und gleichbleibende Capability-Grenzen dürfen in
die Neuplanung eingehen.

\subsection{Vergleichende Trade-off-Diskussion}

Das Episodengedächtnis verbessert vor allem Langzeitpersonalisation und
Promptökonomie, vergrößert aber Datenschutz- und Poisoning-Risiken dauerhaft.
Die Replan-Schleife erhöht die Kosten pro problematischem Turn, begrenzt dafür
unmittelbar Fehlwirkungen und macht Toolfehler beobachtbar. Priorität hat daher
Vorschlag~2: Er beseitigt die sicherheitsrelevante Post-hoc-Kontrolle und schafft
den typisierten Outcome, den Vorschlag~1 später speichern kann. Danach sollte
das Gedächtnis zunächst ohne automatisch generierte Reflexionen und nur für
validierte Aufgabenereignisse eingeführt werden. Beide Erweiterungen sind erst
nach Vergleich gegen den bestehenden Workflow sinnvoll, relevante Kriterien
sind Planerfüllung, Toolpräzision, Revisionsrate, Latenz, Tokenkosten und
mehrsitziger Lernfortschritt. Ohne solche Messungen wäre der zusätzliche
Architekturaufwand nicht begründet.

\section{LLM, RAG und Tool Use (Kap. 03)}
\label{sec:llm-rag-tool-use}

\subsection{LLM- und RAG-Entscheidung}

Als Standardmodell wird \texttt{gpt-5.6-luna} verwendet, alternativ kann ein OpenAI-kompatibler HAW-Endpunkt
verwendet werden. Für den Lernchat ist diese Wahl plausibel als Kosten- und
Latenzkompromiss, weil im Normalfall drei Spezialistenaufrufe und bei der
voreingestellten zweiten Koordinationsrunde sechs entstehen. Konfigurierbar sind
bis zu zwölf Aufrufe, jeweils zuzüglich der internen Toolrunden. Benötigt werden
JSON-Modus und Function Calling, nicht
maximale Modellgröße. Modellabhängige Kompatibilität ist sichtbar: Reasoning-Modelle erhalten keinen
Temperaturparameter. Für GPT-5-Function-Calls setzt der Adapter im
Chat-Completions-Endpunkt \texttt{reasoning\_effort=none}. Ein empirischer Qualitätsvergleich verschiedener Modelle
fehlt. Ihre Eignung ist daher nicht empirisch belegt.

RAG ist fachlich begründet, weil private Lernmaterialien weder zuverlässig
im Modellwissen liegen noch vollständig in jeden Prompt passen. Lewis et al.
beschreiben die Verbindung parametrischen Modellwissens mit explizitem,
nichtparametrischem Speicher \parencite{lewisEtAl2020rag}. StudyMummy indexiert
in einer Offline-Phase Dokumentchunks mit dem konfigurierten Embedding-Modell
(Default \texttt{text-embedding-3-small}) in \texttt{pgvector}. Online folgen
Retrieval, Prompt-Augmentation und Antwortgenerierung: Pro Frage werden die drei
ähnlichsten, nutzerisolierten Chunks geladen. Der Nutzen ist
projektspezifische Grundierung, keine Wahrheitsgarantie. Grenzen sind starres
Zeichen-Chunking, fehlende Ähnlichkeitsschwelle und Quellenangabe, kein
Reranking und stilles Degradieren bei Fehlern. Auch korrekter Kontext schließt
Halluzinationen nicht aus. Embedding und Retrieval erhöhen zudem Latenz und
Kosten.

\subsection{Tools und Agentic Pattern}

Die Registry enthält sechs Werkzeuge. \Cref{lst:tools} zeigt die für Function
Calling relevanten, vereinfachten Schnittstellen. Vier davon liegen abhängig
von der geplanten Aktion im effektiven Chat-Scope.

\begin{lstlisting}[language=Python,float=htbp,caption={Registrierte Werkzeugsignaturen und effektiver Chat-Scope.},label={lst:tools}]
# vom Chat-Workflow freigebbar
evaluate_answer(task_id: str, user_answer: str,
                 expected_concept: str, help_level: int = 1) -> dict
update_learning_profile(tag: str, score: float,
                        error_pattern: str = "") -> dict
award_coins_and_exp(amount: int, reason: str,
                    task_id: str | None = None) -> dict
add_calendar_note(title: str, content: str,
                  start_time: str, end_time: str) -> dict

# registriert, aber durch SAFE_TOOL_POLICY im Chat nicht erreichbar
generate_quiz_questions(topic: str, num_questions: int = 5,
                         difficulty: int = 2) -> dict
create_cheatsheet(user_id: str, session_id: str,
                  topics: list[str]) -> dict
\end{lstlisting}

Das äußere Muster ist eine begrenzte
\emph{Plan--Act--Observe--Review--Replan}-Schleife. Der Tutor alterniert intern
Function Calls und Beobachtungen. Der Reviewer kann den Entwurf akzeptieren,
toolfrei revidieren lassen oder Beobachtung und Kritik an den Planner
zurückgeben. Dies überträgt ReActs Kopplung von Reasoning und Handeln auf den
typisierten MAS-Kontrollfluss \parencite{yaoEtAl2023react}, ohne interne
Gedankengänge zu veröffentlichen. \texttt{tool\_calls} enthält nur erfolgreiche
Funktionen. Alle Versuche werden separat als \texttt{ToolObservation} mit
Status und gekürzter Resultatvorschau an Reviewer und API gemeldet. Plan,
sanitisierte Nachrichtenmetadaten, lokale Agentenzustände und
Koordinationsrunden sind beobachtbar, werden aber weiterhin nicht als
vollständige Episode persistiert.

\section{Multi-Agent-Entscheidung (Kap. 04)}
\label{sec:multi-agent}

StudyMummy ist in der untersuchten Implementierung als \emph{hierarchisches,
nachrichtengetriebenes Multi-Agent-System} einzuordnen. Entscheidend ist nicht
die Zahl der Prompts, sondern ausführbare Agentenmerkmale: Planning, Tutor und
Reviewer implementieren denselben \texttt{MASAgent}-Vertrag, besitzen je eine
Identität, ein eigenes Ziel, Capabilities und einen getrennten
\texttt{AgentLocalState}. Dieser Zustand zählt empfangene und gesendete
Nachrichten sowie Entscheidungen und hält eine kleine lokale Erinnerung für
den aktuellen Turn. Alle drei nehmen einen Nachrichtentyp wahr, entscheiden
innerhalb ihrer Rolle und erzeugen selbst Folgebotschaften.

Kooperation erfolgt über einen run-lokalen \texttt{MASMessageBus} mit
Performatives wie \texttt{request}, \texttt{delegate}, \texttt{propose},
\texttt{critique} und \texttt{accept}. Ein gemeinsames
\texttt{AgentBlackboard} repräsentiert die veränderliche Umgebung aus Plan,
Entwurf, Review und Toolbeobachtungen. Der normale Pfad lautet Planner--Tutor--Reviewer.
Bei Dissens routet der Reviewer autonom eine \texttt{revision\_request} an den
Tutor oder eine \texttt{replan\_request} an den Planner. Dadurch verändert er
den Kontrollfluss tatsächlich, statt lediglich einen Text zu ersetzen. Der
deterministische Orchestrator ist kein vierter Agent: Er dispatcht FIFO,
erfasst Laufzeiten und beendet nach höchstens der konfigurierten Rundenzahl.
Perception, Memory und RAG bleiben Umgebungsdienste, nicht künstlich
umbenannte Agenten.

\begin{center}
\begin{minipage}{\linewidth}
  \centering
  \begin{tikzpicture}[
    >=Latex,
    box/.style={draw,rounded corners,align=center,minimum height=8mm,text width=25mm,font=\small},
    role/.style={box,fill=hawblue!8},
    service/.style={box,fill=gray!10}
  ]
    \node[box] (ui) at (-4.8,1.7) {Lernende Person\\Angular};
    \node[box] (api) at (-1.6,1.7) {FastAPI\\\texttt{/agent/chat}};
    \node[service] (env) at (1.6,1.7) {Perception\\Session + RAG};
    \node[service] (memory) at (4.8,1.7) {Memory\\Chatlog + DB};
    \node[service,text width=31mm] (bus) at (0,0) {MAS Runtime\\Message Bus + Blackboard};
    \node[role] (planner) at (-4.2,-1.8) {Planner Agent\\Plan/Replan};
    \node[role] (tutor) at (0,-1.8) {Tutor Agent\\Act/Revise};
    \node[role] (reviewer) at (4.2,-1.8) {Reviewer Agent\\Accept/Critique};
    \node[service] (tools) at (-2.2,-3.6) {Tool Registry\\Side Effects};
    \node[service,text width=34mm] (llm) at (2.2,-3.6) {gemeinsamer OpenAI-kompatibler LLM-Adapter};

    \draw[->,semithick] (ui) -- (api);
    \draw[->,semithick] (api) -- (env);
    \draw[->,semithick] (env) -- (bus);
    \draw[->,semithick] (api) -- (memory);
    \draw[<->,semithick] (bus) -- (planner);
    \draw[<->,semithick] (bus) -- (tutor);
    \draw[<->,semithick] (bus) -- (reviewer);
    \draw[->,semithick] (tutor) -- (tools);
    \draw[dashed,->] (llm) -- (planner);
    \draw[dashed,->] (llm) -- (tutor);
    \draw[dashed,->] (llm) -- (reviewer);
  \end{tikzpicture}
  \captionsetup{hypcap=false}
  \captionof{figure}{Selbst erstellte Ist-Architektur des hierarchischen MAS.
    Rückpfeile am Bus stehen für Revision und Neuplanung. Volle
    Nachrichten-Payloads werden nicht an die API serialisiert.}
  \label{fig:mas-architektur}
\end{minipage}
\end{center}

\Cref{fig:mas-architektur} zeigt zugleich die Abgrenzung: Das MAS ist homogen,
seriell und in-process. Die Agenten teilen sich einen Modelladapter, besitzen
keine eigenen Prozesse und ihr lokaler Zustand endet mit dem Turn. Es gibt
weder Parallelität, freie Aushandlung noch horizontale Fehlerisolation. Ein
Single-Agent-Entwurf wäre günstiger, würde jedoch die getrennten Ziele,
Capability-Grenzen und den überprüfbaren Dissenspfad verlieren. Verteilte
Agentendienste wären umgekehrt ohne Parallelitäts- oder Verfügbarkeitsziel
unnötig komplex. Die gewählte Mitte liefert echte, aber bewusst begrenzte
Kooperation und bleibt für einen Tutor-Turn testbar.

\section{Engineering und Testing (Kap. 05)}
\label{sec:engineering-testing}

\subsection{Framework- und Implementierungsentscheidungen}

Die MAS-Laufzeit ist bewusst ohne zusätzliches Agentenframework implementiert.
\texttt{MASAgent} definiert den gemeinsamen Agentenvertrag,
\texttt{MASMessageBus} eine FIFO-Mailbox mit sanitisiertem Kommunikations-Trace
und \texttt{AgentBlackboard} den gemeinsamen Laufzustand. Der Orchestrator
kennt Agentenregister, Nachrichtenwarteschlange und Abbruchgrenze, aber nicht
die fachliche Entscheidung: Das Routing nach Kritik erzeugt der Reviewer.
Pydantic validiert Pläne, Reviews, Nachrichten, Toolstatus, lokale Zustände und
API-Antworten. Sämtliche Laufobjekte werden je Request erzeugt. Dadurch teilt
der globale FastAPI-Orchestrator keinen veränderlichen Agentenzustand zwischen
parallelen Nutzern.

Diese Eigenimplementierung hält den kleinen Zustandsgraphen und die
Capability-Policy direkt prüfbar. Ein graphorientiertes Framework wäre bei
persistenten Checkpoints, vielen Verzweigungen oder Human-in-the-Loop-Pausen
neu abzuwägen. Gegenwärtig würde es zusätzliche Abhängigkeit und Migration
erzeugen, ohne eine benötigte Funktion zu ersetzen. Der Preis des eigenen
Runtimes bleibt Wartungsaufwand: Zustandsmigration, Wiederaufnahme und
verteilte Ausführung müssten selbst entwickelt werden. Das System ist
asynchron programmiert, dispatcht die Agenten aber absichtlich seriell.

Der OpenAI-kompatible Adapter kapselt JSON-Modus und Function Calling. Die
Tool-Registry trennt Schema und Implementierung, während
\texttt{SAFE\_TOOL\_POLICY} Rechte aktionsabhängig reduziert. Toolresultate
werden als Erfolg, Fehler, Blockade oder Argumentfehler typisiert. Der Reviewer
kann auf dieser Basis Textrevision oder Neuplanung auslösen. Ein Limit von eins
bis vier, standardmäßig zwei Koordinationsrunden verhindert Endlosschleifen.
Weiterhin fehlen persistente Checkpoints, Retries mit Backoff und ein
transaktionales Zurückrollen bereits ausgeführter Tools.

\subsection{Testfälle und Beobachtungen}

Alle 64 Tests bestanden. Die einzige Warnung betrifft die künftige
Umstellung des Starlette-Testclients. Zusätzlich bestanden zehn Angular-Tests
und der Angular-Produktionsbuild. Beide Docker-Compose-Profile wurden
erfolgreich aufgelöst. Die drei in \Cref{tab:tests} ausgewählten Tests
adressieren gerade die Merkmale, die eine Rollen-Pipeline nicht belegen könnte.

Ein ergänzender Live-End-to-End-Lauf mit dem konfigurierten OpenAI-Modell
verarbeitete ein zweiseitiges Test-PDF zu vier Aufgaben, einem Quiz mit sechs
Fragen und einem Merkblatt. Der Hinweis-Turn beteiligte Planner, Tutor und
Reviewer über sechs Nachrichten und zwei Runden. Im Bewertungsturn wurde
\texttt{evaluate\_answer} mit Status \texttt{succeeded} protokolliert. Eine
dokumentgestützte Multiple-Choice-Übung wurde ebenfalls erzeugt. Der Lauf ist
ein Integrationsnachweis, aber keine Nutzerstudie.

\begin{table}[htbp]
  \centering
  \footnotesize
  \begin{tabularx}{\linewidth}{@{}p{0.25\linewidth}XX@{}}
    \toprule
    Testfall & Prüfkriterium & Beobachtung \\
    \midrule
    Nachrichten und lokaler Zustand & Erwartet werden die typisierte Folge
    \texttt{plan\_request}, \texttt{execute\_plan},
    \texttt{review\_request}, \texttt{final\_response} sowie je eine lokale
    Entscheidung der drei Agenten. & Test bestanden. Die Empfängerfolge war
    Planner, Tutor, Reviewer, Coordinator. \\
    Reviewer-Revision & Ein Entwurf verrät die Lösung. Der Reviewer muss Kritik
    an den Tutor senden. Die toolfreie Revision wird erneut geprüft. & Test
    bestanden. Es gab zwei Runden, zwei Tutor- und zwei Reviewer-Entscheidungen. Die
    textbezogene Assertion prüft den Beginn der ausgelieferten Frage. \\
    Neuplanung nach Toolfehler & Eine fehlgeschlagene
    \texttt{evaluate\_answer}-Beobachtung muss den Reviewer zu
    \texttt{replan\_request} und den Planner zu einem neuen
    \texttt{clarify}-Plan veranlassen. & Test bestanden. Der zweite Plan und
    die Fehlerbeobachtung waren im Laufzustand nachweisbar. \\
    \bottomrule
  \end{tabularx}
  \caption{Ausgewählte, tatsächlich ausgeführte MAS-Tests.}
  \label{tab:tests}
\end{table}

Die markerbasierten semantischen Assertions des kontrollierten Experiments
nutzen vorgegebene Modellantworten, aber den produktiven Tool- und
MAS-Kontrollfluss. Die
ReAct-Variante erreichte bei drei konstruierten Läufen Score und Tool-Abdeckung
von 1,00. Die handgefertigte Baseline erreichte 0,27 beziehungsweise 0,00. Ein zusätzlicher
MAS-Lauf erzeugte reproduzierbar \texttt{plan\_request} $\rightarrow$
\texttt{execute\_plan} $\rightarrow$ \texttt{review\_request} $\rightarrow$
\texttt{revision\_request} $\rightarrow$ \texttt{review\_request}
$\rightarrow$ \texttt{final\_response} und endete nach zwei Runden mit einer
sokratischen Frage. Dies belegt
Nachrichtenrouting und Rückkopplung, nicht fachliche Modellqualität.

Aussagen zu Tutorqualität und Lernfortschritt benötigen deshalb Live-LLM- und
Nutzerstudien. Die Angular-Suite prüft UI-Komponenten. Jede neue Tutorantwort
zeigt außerdem eine aufklappbare Zusammenfassung aus beteiligten
Agenten, Rundenzahl, Nachrichtentypen und Toolstatus. Vollständige Payloads
bleiben auf die interne Laufzeit beschränkt.

\section{Safety-Analyse (Kap. 06)}
\label{sec:safety}

\subsection{Angriffsvektoren und Guardrails}

\textbf{Indirekte Prompt-Injection aus Lernmaterialien.} Ein Angreifer kann in
einem hochgeladenen PDF Anweisungen platzieren, die beim semantischen Abruf in
den Tutor-Prompt gelangen. Genau diese Vermischung von Daten und Instruktionen
ermöglicht indirekte Prompt-Injection bis hin zur Manipulation nachgelagerter
API-Aufrufe \parencite{greshakeEtAl2023}. StudyMummy filtert einige direkte
Muster in Nutzernachrichten, isoliert RAG nach Nutzer und Dokument und markiert
Dokumenttext im Systemprompt als nicht vertrauenswürdig. Der Reviewer soll
befolgte Dokumentanweisungen erkennen. Diese Guardrails reduzieren das Risiko,
sind aber nicht vollständig: Chunks und \texttt{extra\_context} werden nicht
semantisch geprüft, und Review findet erst nach Toolausführung statt.

Passend wären mehrere Grenzen: Dokumente und Frontend-Kontext immer als
unvertrauenswürdige, zitierbare Datensätze serialisieren, Retrieval-Ergebnisse
mit Provenienz versehen, geplante Toolargumente vor Ausführung gegen eine
deterministische Policy prüfen, Dokumentkontext darf niemals neue Tools oder
Identitäten bestimmen. Ergänzend sollten Injection-Testkorpora, Rate Limits
und Quarantäne bei verdächtigen Chunks eingesetzt werden. Ein weiterer
LLM-Filter allein wäre wegen korrelierter Fehlentscheidungen kein verlässlicher
Sicherheitsanker.

\textbf{Toolmissbrauch und ungewollte Seiteneffekte.} Direkte Aufforderungen
könnten Coins erzeugen oder ohne ausreichende Bestätigung Termine anlegen.
Vorhanden sind aktionsabhängige Toolrechte, eine zweite Scope-Prüfung im
LLM-Adapter, gebundene Nutzeridentität sowie beim Belohnungstool Betragsgrenze,
Eigentumsprüfung, Zeilensperre und Doppelbelohnungsschutz. Die einfache Suche
nach Wörtern wie \enquote{coin} ist umgehbar und kann legitime Nachrichten
blockieren. Beim Kalender fehlen Nutzerbestätigung, Idempotenzschlüssel und die
fachliche Prüfung, dass Ende nach Beginn liegt. Da der Reviewer erst nach der
Ausführung arbeitet, kann er diese Wirkung nicht verhindern.

Ein Pre-Execution-Monitor sollte daher Tool, Argumente, Identität,
Wertebereiche und erwartete Aktion prüfen. Schreibende Aufrufe werden zunächst
als \texttt{pending} protokolliert, Kalendertermine werden erst nach
Bestätigung committet. Für Belohnungen reichen wegen geringer Reversibilität
enge automatisierte Grenzen und Auditlog. Verbleibende Risiken sind
LLM-Nichtdeterminismus, Promptvarianten außerhalb der Filter, vergiftetes
Langzeitgedächtnis, Kostenmissbrauch durch große Uploads sowie personenbezogene
Inhalte in Logs und Vektorspeicher.

\subsection{Human-in-the-Loop}

Vollständige Freigabe jedes Tutor-Turns würde die Lerninteraktion unbrauchbar
verlangsamen. Geeignet ist ein risikobasiertes Human-in-the-Loop: Lesen,
Retrieval, Antwortbewertung und Textentwürfe laufen autonom, ein sichtbarer
Bestätigungsdialog zeigt bei persistenten Seiteneffekten Toolname, Argumente,
Quelle und erwartete Wirkung. Die lernende Person bestätigt Kalendertermine
und kann gespeicherte Erinnerungen löschen. Lehrende oder Betreiber prüfen
stichprobenartig abgelehnte Antworten und Injection-Traces. HITL begrenzt den
Schadensradius, ersetzt aber weder Rechteprüfung noch Monitoring: Menschen
können unter Zeitdruck routinemäßig bestätigen und sehen versteckte
Kontextmanipulation nicht. Deshalb muss die technische Policy auch bei
menschlicher Freigabe unverändert gelten.

\section*{Synthese und Fazit}
\addcontentsline{toc}{section}{Synthese und Fazit}

StudyMummy koppelt die vier kognitiven Schichten in einem hierarchischen,
nachrichtengetriebenen Multi-Agent-System: Perception bildet den Lernzustand ab,
Sitzungen, Chatlogs und Vektoren tragen Memory, der Planner plant und der Tutor handelt.
Planner, Tutor und Reviewer besitzen getrennte Ziele, Capabilities und lokale
Zustände. Über typisierte Nachrichten führen Toolbeobachtungen und Reviewkritik
innerhalb eines Turns zu Revision oder Neuplanung. Der Regelkreis ist damit bis
zur konfigurierten Rundengrenze geschlossen.

Begrenzt bleibt das MAS durch serielle In-Process-Ausführung, einen gemeinsamen
LLM-Adapter und flüchtige Agentenzustände. Lernprofil und frühere Outcomes
steuern den Plan nicht. Schreibende Tools wirken vor dem Review und lassen sich
dort nicht zurückrollen. Vorrang haben deshalb ein transaktionaler Prüf- und
Commitpfad sowie danach ein Gedächtnis für validierte Outcomes. Die Tests
belegen Routing, Rückkopplung und Verträge, aber weder Tutorqualität noch
Lernwirksamkeit. Dafür fehlen Live-LLM- und Nutzerstudien.

\section*{Erklärung zur Nutzung von KI-Werkzeugen}
\addcontentsline{toc}{section}{Erklärung zur Nutzung von KI-Werkzeugen}

Für diese Arbeit habe ich OpenAI ChatGPT 5.6 zur Auswertung der Aufgabenstellung, zur Hilfe beim Programmieren und
Anpassen des Codes eingesetzt. Zusätzlich half mir ChatGPT bei der Strukturierung, Formulierungsfragen, sowie der Literaturrecherche. Die Verantwortung für
Auswahl, Bewertung und eingereichte Fassung trage ich selbst.

\printbibliography[heading=bibintoc,title={Literaturverzeichnis}]

\end{document}
